% ============================================================================
% preamble.tex — 物理习题集
% ============================================================================

% --- 中文支持 ---
\usepackage[fontset=fandol]{ctex}
\usepackage{fontspec}

% --- 页面布局 ---
\usepackage[a4paper, top=2.5cm, bottom=2.5cm, left=2.5cm, right=2.5cm]{geometry}

% --- 数学 ---
\usepackage{amsmath, amssymb, amsthm}
\usepackage{mathrsfs}
\usepackage{mathtools}
\usepackage{bm}
\usepackage{physics}

% --- 图表 ---
\usepackage{graphicx}
\usepackage{float}
\usepackage{booktabs}
\usepackage{longtable}
\usepackage{tabularx}

% --- 绘图 ---
\usepackage{tikz}
\usetikzlibrary{
  positioning, calc, arrows.meta,
  decorations.pathreplacing, decorations.markings,
  patterns, angles, quotes,
}
\usepackage[american]{circuitikz}

% --- 代码 ---
\usepackage{listings}
\lstset{
  language=Python,
  basicstyle=\ttfamily\small,
  keywordstyle=\color{blue!70!black},
  commentstyle=\color{green!50!black},
  stringstyle=\color{red!60!black},
  numbers=left, numberstyle=\tiny\color{gray},
  frame=single, breaklines=true,
}

% --- 单位 ---
\usepackage{siunitx}
\sisetup{per-mode=symbol}

% --- 超链接 ---
\usepackage{hyperref}
\hypersetup{
  colorlinks=true,
  linkcolor=blue!60!black,
  citecolor=green!40!black,
  urlcolor=cyan!50!black,
}
\usepackage[nameinlink]{cleveref}

% --- 索引 ---
\usepackage{makeidx}
\makeindex

% --- 颜色与框 ---
\usepackage{xcolor}
\usepackage[most]{tcolorbox}
\tcbuselibrary{breakable, skins, listings}

% ============================================================================
% 习题环境
% ============================================================================

\newcounter{problem}[chapter]
\renewcommand{\theproblem}{\thechapter.\arabic{problem}}

% 普通习题（蓝色框）
\newtcolorbox{problem}[1][]{
  enhanced, breakable, arc=3pt,
  before skip=12pt, after skip=12pt,
  colback=blue!3, colframe=blue!40!black,
  boxrule=0.6pt, left=8pt, right=8pt, top=8pt, bottom=8pt,
  title={\refstepcounter{problem}\bfseries 题 \theproblem},
  fonttitle=\large, coltitle=white, colbacktitle=blue!50!black,
  #1
}

% 经典题（紫色框）
\newtcolorbox{classic}[1][]{
  enhanced, breakable, arc=3pt,
  before skip=12pt, after skip=12pt,
  colback=violet!3, colframe=violet!50!black,
  boxrule=0.6pt, left=8pt, right=8pt, top=8pt, bottom=8pt,
  title={\refstepcounter{problem}\bfseries 经典题 \theproblem},
  fonttitle=\large, coltitle=white, colbacktitle=violet!60!black,
  #1
}

% 竞赛题（红色框）
\newtcolorbox{competition}[1][]{
  enhanced, breakable, arc=3pt,
  before skip=12pt, after skip=12pt,
  colback=red!3, colframe=red!50!black,
  boxrule=0.8pt, left=8pt, right=8pt, top=8pt, bottom=8pt,
  title={\refstepcounter{problem}\bfseries 竞赛题 \theproblem},
  fonttitle=\large, coltitle=white, colbacktitle=red!60!black,
  #1
}

% 大学物理题（绿色框）
\newtcolorbox{university}[1][]{
  enhanced, breakable, arc=3pt,
  before skip=12pt, after skip=12pt,
  colback=green!3, colframe=green!40!black,
  boxrule=0.6pt, left=8pt, right=8pt, top=8pt, bottom=8pt,
  title={\refstepcounter{problem}\bfseries 大学物理 \theproblem},
  fonttitle=\large, coltitle=white, colbacktitle=green!50!black,
  #1
}

% 解答环境
\newenvironment{solution}{%
  \par\noindent\textbf{解：}\par\nopagebreak
}{\par\medskip}

% ============================================================================
% 物理符号
% ============================================================================
\newcommand{\vect}[1]{\vb{#1}}
\newcommand{\uvec}[1]{\vu{#1}}
\newcommand{\magn}[1]{\left|#1\right|}
\newcommand{\avg}[1]{\langle #1 \rangle}
